\documentclass[12pt, letterpaper]{article}

\usepackage{amsmath, amssymb}
\usepackage{physics}
\usepackage[french]{babel}
\usepackage[margin=2.5cm]{geometry}
\usepackage[utf8]{inputenc}
\usepackage[T1]{fontenc}
\usepackage{hyperref}
\usepackage{caption}
\usepackage{booktabs}
\usepackage{multirow}
\usepackage{tcolorbox}

\begin{document}

\title{Devoir 3 : Équation de Schrödinger en dimension 1}
\author{PHQ404}
\date{Date de remise : -}
\maketitle

\section{Objectif}\label{sec:objectif}

\noindent L’objectif de ce travail est de se familiariser avec deux techniques de résolution
de problèmes aux limites en dimension 1 : la méthode du tir et la méthode des éléments finis.
Ces méthodes sont présentées au chapitre 12 des notes de David Sénéchal.

\section{Définition du problème}\label{sec:definition-probleme}

\noindent L’équation de Schrödinger indépendante du temps est considérée pour une particule de masse unitaire dans
un potentiel quadratique :
\begin{equation}
	\label{eq:schrodinger}
	- \frac{1}{2} \ddot{\psi} + \frac{1}{2} x^2 \psi = E \psi.
\end{equation}
Les solutions analytiques sont connues et les énergies propres exactes sont
\begin{equation}
	E_n = n + \tfrac{1}{2}, \qquad n = 0,1,2,\dots
\end{equation}
Ces valeurs serviront à valider vos résultats numériques.

\section{Méthode du tir}\label{sec:methode-du-tir}

\noindent L'équation de Schrödinger doit être résolue numériquement avec la méthode du tir.
On cherche les solutions pour $x \in [-L,L]$ avec $L = 5$, en utilisant $N=200$ points uniformément espacés.
Les conditions initiales utilisées pour l'intégration sont
\begin{equation}
	\psi(-L) = 0, \qquad \dot{\psi}(-L) = 0{,}001.
\end{equation}
La valeur initiale de $\dot{\psi}$ n’est pas critique et n'influence pas les énergies propres,
mais celle-ci devrait vous donner une échelle raisonnable. \\

\noindent Les énergies propres sont définies comme les valeurs de $E$ telles que la solution intégrée satisfait la condition de Dirichlet $\psi(L \mid E) = 0$.
On définit donc la fonction
\begin{equation}
	f(E) = \psi(L \mid E),
\end{equation}
et les énergies propres correspondent aux racines de $f(E)$. \\

\noindent Vous devez :

\begin{itemize}
	\item Implémenter une fonction qui, pour une valeur donnée de $E$, intègre l'équation et retourne $f(E) = \psi(x=L \mid E)$.

	\item Implémenter une procédure de cadrage de racine:
	      étant donnée une valeur initiale $E_0$, trouver $E_1$ tel que
	      \begin{equation}
		      f(E_0)\,f(E_1) < 0.
	      \end{equation}
	      Pour ce faire, augmenter $E$ par pas de $\Delta E$ jusqu’à détecter un changement de signe,
	      ou atteindre un nombre maximal d’itérations (20 par exemple).

	\item Utiliser ensuite l'algorithme de Brent pour trouver la racine dans l'intervalle ainsi obtenue.
	      L’implémentation de \texttt{SciPy} peut être utilisée pour l’algorithme de Brent.

	\item Déterminer numériquement les six plus basses énergies propres.

	\item Représenter sur un même graphique les six fonctions d’onde correspondantes,
	      chacune étant décalée verticalement par la valeur de son énergie.
\end{itemize}

\begin{figure}[h!!]
	\centering
	\includegraphics[scale=0.5]{../images/exemple}
	\caption{Exemple de figure pour les états propres de basses énergies
		de l’équation~\ref{eq:schrodinger}.}
	\label{fig:exemple}
\end{figure}

\section{Méthode des éléments finis}\label{sec:methode-des-elements-finis}

\noindent Dans cette partie, la méthode des éléments finis doit être utilisée pour la résolution de l'équation de Schrödinger.
Un fichier Python vous est fourni et contient une classe implémentant les éléments nécessaires. \\

\noindent Vous devez :

\begin{itemize}
	\item Utiliser cette classe pour résoudre numériquement l’équation~\ref{eq:schrodinger}.
	\item Utiliser une grille de $N=500$ points.
	\item Prendre $L=6$.
	\item Produire le même type de graphique que dans la partie précédente.
\end{itemize}

\noindent Les fonctions fournies ne calculent pas les valeurs aux frontières.
Celles-ci ne sont pas nécessaires et peuvent ajouter du bruit numérique aux solutions.

\begin{tcolorbox}
	Toutes les fonctions qui doivent être implémentées sont déjà définies dans les fichiers
	et retournent une exception \texttt{NotImplementedError}.
\end{tcolorbox}

\section{Comparaison et vérification}\label{sec:verification}

Comment est-ce que les deux méthodes se comparent?
Comment peut-on bien voir les différences et analyser le taux d'erreur entre celles-ci?
Quelle méthode converge le plus rapidement?
Quelle sont les compromis en termes de complexité algorithmique?
Comment est-ce que ces méthodes se comparent à la solution analytique?
La fonction \texttt{scipy.special.eval\_hermite} pourrait vous être utile.
\begin{tcolorbox}[]
	Des exemples de vérification sont fournis, mais il vous appartient d’ajouter un nombre suffisant de tests unitaires pertinents ainsi que des tests de validation scientifique afin d’obtenir l’ensemble des points.
\end{tcolorbox}

\section{Barème d'évaluation}\label{sec:bareme-evaluation}

\begin{tabular}{l l c}
	\toprule
	\textbf{Catégorie} & \textbf{Critère d’évaluation}              & \textbf{Pondération} \\
	\midrule
	\multirow{6}{*}{Code (60\%)}
	                   & Qualité du code (pylint)                   & 10\%                 \\
	                   & Couverture de tests (codecov)              & 10\%                 \\
	                   & Tests publics (pytest)                     & 10\%                 \\
	                   & Tests cachés (pytest)                      & 20\%                 \\
	                   & Documentation, docstrings et typage        & 5\%                  \\
	                   & README, reproductibilité et historique Git & 5\%                  \\
	\midrule
	\multirow{9}{*}{Rapport (40\%)}
	                   & Résumé                                     & 3\%                  \\
	                   & Introduction                               & 3\%                  \\
	                   & Théorie                                    & 6\%                  \\
	                   & Résultats                                  & 6\%                  \\
	                   & Discussion                                 & 6\%                  \\
	                   & Conclusion                                 & 3\%                  \\
	                   & Références                                 & 3\%                  \\
	                   & Vérification des résults                   & 5\%                  \\
	                   & Présentation                               & 5\%                  \\
	\bottomrule
\end{tabular}


\end{document}
